\abstract{РЕФЕРАТ}

Объем работы равен \formbytotal{lastpage}{страниц}{е}{ам}{ам}. Работа содержит \formbytotal{figurecnt}{иллюстраци}{ю}{и}{й}, \formbytotal{tablecnt}{таблиц}{у}{ы}{}, \arabic{bibcount} библиографических источников и \formbytotal{числоПлакатов}{лист}{}{а}{ов} графического материала. Количество приложений – 1. Фрагменты исходного кода представлены в приложении А.

Перечень ключевых слов: продукционная система Маркова, моделирование, алгоритм, правила замены, логирование, статистика, Python, Tkinter, CustomTkinter, графический интерфейс, пошаговое выполнение, анимация, визуализация, строка, правило, левый и правый символ, конфигурация.

Объектом разработки является программная система на языке Python для симуляции работы продукционной системы Маркова с графическим интерфейсом и возможностью визуализации пошагового выполнения алгоритмов.

Целью работы является создание программной реализации продукционной системы Маркова для демонстрации принципов работы формальных алгоритмических систем, изучения применения правил замены и визуального анализа выполнения алгоритмов.

В процессе разработки были реализованы основные компоненты системы: правила Маркова (левая и правая части), строка для преобразования, механизм применения правил, подсчет статистики выполнения (количество шагов, частота применения правил, длина строки на каждом шаге), журнал логирования всех шагов, а также графический интерфейс для визуализации процесса с использованием Tkinter и CustomTkinter.

Разработанное приложение позволяет загружать правила из текстовых файлов, вводить правила вручную, выполнять алгоритм пошагово с анимацией, сохранять журнал работы и отображать статистику. Программная реализация успешно демонстрирует работу продукционной системы Маркова на различных примерах алгоритмов.

\selectlanguage{english}
\abstract{ABSTRACT}

The volume of work is \formbytotal{lastpage}{page}{}{s}{s}. The work contains \formbytotal{figurecnt}{illustration}{}{s}{s}, \formbytotal{tablecnt}{table}{}{s}{s}, \arabic{bibcount} bibliographic sources and \formbytotal{числоПлакатов}{sheet}{}{s}{s} of graphic material. The number of applications is 2. The graphic material is presented in annex A. Fragments of the source code are presented in annex B.

List of keywords: Markov production system, modeling, algorithm, replacement rules, logging, statistics, Python, Tkinter, CustomTkinter, graphical user interface, step-by-step execution, animation, visualization, string, rule, left and right parts, configuration.

The object of development is a software system implemented in Python for simulating a Markov production system with a graphical interface and the ability to visualize step-by-step algorithm execution.

The purpose of the work is to create a software implementation of a Markov production system to demonstrate the principles of formal algorithmic systems, study the application of replacement rules, and perform visual analysis of algorithm execution.

During development, the main components of the system were implemented: Markov rules (left and right parts), the working string, the rule application mechanism, execution statistics (number of steps, frequency of rule application, string length at each step), logging of all steps, and a graphical interface for visualizing the process using Tkinter and CustomTkinter.

The developed application allows loading rules from text files, manual rule input, step-by-step execution with animation, saving execution logs, and displaying statistics. The software implementation successfully demonstrates the operation of a Markov production system on various algorithm examples.
\selectlanguage{russian}